Here, we observe the difference of paper rankings to further examine the evolution of highly cited papers in detail. 
By performing year-by-year comparisons using Kendall's rank correlation, a higher correlation coefficient implies less variation in rankings, i.e., greater stability. 
The subjects were the highly cited papers from all accumulated generations. For each year from 1947 to 2020, we calculated the correlation coefficient relative to the previous year. Figure \ref{fig:gr101} shows the correlation coefficients for each year. The values fluctuated between 0.85 and 0.99, with no discernible trend over the years. 

Furthermore, to examine how the cumulative effect of citations manifests, the figure also displays the correlation coefficients calculated using cumulative citations within $n$ years after publication. 
According to these results, as the citation window (impact years) increases, the correlation coefficient tends to decrease, indicating a tendency for rankings to become unstable. This suggests that the cumulative effect of citations is likely a factor driving fluctuations in paper rankings.
